% !PRINTED 217
\weberpage{217}

%<<CONT
and the relations
\weq{6}{\Theta\chi = \chi^2\psi\Theta^3, \quad \Theta^2\chi = \chi\psi, \quad \Theta^3\chi = \chi^2\Theta \\
\psi\chi = \chi\psi\Theta^2, \quad \psi\chi^2 = \chi^2\Theta^2.}
These relations, in combination with (2), show that the
substitutions
\weq{7}{\chi^{\lambda}\psi^{\mu}\Theta^{\nu}, \quad \lambda = 0, 1, 2; \quad \mu = 0, 1; \quad \nu = 0, 1, 2, 3}
do in fact form a group, because with their aid every
substitution of the form
\[
\chi^{\lambda}\psi^{\mu}\Theta^{\nu}\chi, \quad \chi^{\lambda}\psi^{\mu}\Theta^{\nu}\psi,
\]
and consequently, by repetition, every substitution
\[
\chi^{\lambda}\psi^{\mu}\Theta^{\nu}\,\chi^{\lambda'}\psi^{\mu'}\Theta^{\nu'}
\]
can be brought into the form (7). The group can be represented
explicitly in the form
\weq{8}{i^{\nu}x, \quad \frac{i^{\nu}}{x}, \quad i^{\nu}\,\frac{x - i^{\nu'}}{x + i^{\nu'}}, \quad \nu, \nu' = 0, 1, 2, 3}
and is of the 24th degree. It is the
\gesperrt{octahedral group}.

The group has a normal divisor of the 12th degree, which can
be represented in the form $\chi^{\lambda}\psi^{\mu}\Theta^{2\nu}$ and which is a
tetrahedral group.

The representation becomes in a certain respect simpler if,
in place of $\psi$, one introduces a new element
\weq{9}{\omega = \psi\Theta = \frac{-i}{x}}
which, just like $\psi$, is of the second order. Then
the group can be represented by
\[
\chi^{\lambda}\omega^{\mu}\Theta^{\nu},
\]
and in place of the relations (2) and (6) the following appear:
\[
\omega\chi = \chi^2\omega, \quad \omega\chi^2 = \chi\omega
\]
\weq{10}{\Theta\chi = \chi^2\omega\Theta^2, \; \Theta^2\chi = \chi\omega\Theta^3, \; \Theta^3\chi = \chi^2\Theta, \; \Theta^2\chi^2 = \chi^2\omega\Theta \\
\Theta\omega = \omega\Theta^3, \quad \Theta^2\omega = \omega\Theta^2, \quad \Theta^3\omega = \omega\Theta.}
But all these relations follow as consequences
from the four:
\weq{11}{\omega\chi = \chi^2\omega, \quad \Theta\omega = \omega\Theta^3, \quad \Theta\chi = \chi^2\omega\Theta^2, \quad \Theta^2\chi = \chi\omega\Theta^3,}
in combination with the formulae expressing the degrees:
\weq{12}{\chi^3 = 1, \quad \omega^2 = 1, \quad \Theta^4 = 1.}
